Across sound studies, human–computer interaction, and design research, there is broad recognition that technology shapes perception, but little work connects this understanding to systematic methods for transforming the tools themselves. Post-phenomenology reveals how technologies mediate experience, yet rarely considers how the specifics of cultural epistemologies might intervene in those mediations. Meanwhile, critical race and discriminatory design discourses expose how sound and computation encode power, but often remain at the level of interpretation and critique rather than material modification.

Recent scholarship makes clear that design paradigms embed Western assumptions about temporality, identity, and interaction into computational logic while presenting these assumptions as universal. Frameworks like Mechanics, Dynamics, and Aesthetics (MDA), which is “probably the most widely accepted and practically employed approach to game design” \parencite{Walk_2017}, encodes specific values of agency and player behavior into procedural systems. Platform studies show that tools like Unreal Engine operate as “technical and political archives” that reproduce cultural hierarchies through their very architectures \parencite{Malazita_2024}. Likewise, generative AI exemplifies how “AI universalism creates narrow human subjects” by reducing cultural complexity to extractable data \parencite{Klein_2025}. In digital audio, Jonathan Sterne reminds us that “encoded in every mp3 are whole worlds of possible and impossible sound and whole histories of sonic practice,” where compression algorithms silently decide what sounds matter based on “an imagined, ideal auditor” \parencite[p.2]{Sterne_2012}. Such defaults naturalize particular ways of hearing, seeing, and knowing while rendering others technically illegible.

What remains missing is a framework that bridges critique and creation as a methodology for reading and rewriting the logics within our digital instruments. While scholars have mapped how cultural bias operates in technological systems, practitioners still lack systematic approaches for engaging those biases through practice. This research responds to that absence by developing \emph{Speculocultural Technopoiesis} as a mode of research-creation that unites cultural theory with technical modification. It positions sonic and procedural making as a continuous act of epistemological translation, transforming speculative, culturally grounded imagination into new computational and auditory forms through collaboration, autoethnography, and acts of "re-tuning."